% © 2026 Ing. David Jaroš — CC BY-NC-ND 4.0
% UBT-AI-PROVENANCE-BEGIN
% schema: ubt-ai-provenance/v1
% tier: C_working
% ai_assistance: disclosed
% human_review: risk-based
% editorial_responsibility: Ing. David Jaroš
% policy: ../../../AI_PROVENANCE.md
% notice: Working material; exhaustive human review is not claimed.
% UBT-AI-PROVENANCE-END

\documentclass[12pt,a4paper]{article}
\usepackage[a4paper,margin=2.5cm]{geometry}
\usepackage{amsmath,amssymb,amsthm,mathtools,booktabs,longtable}
\usepackage[most]{tcolorbox}
\usepackage{hyperref}
\usepackage{xcolor}
\usepackage{microtype}
\hypersetup{colorlinks=true,linkcolor=blue!60!black,urlcolor=blue!60!black,citecolor=green!50!black}
\newtheorem{theorem}{Theorem}[section]
\newtheorem{lemma}[theorem]{Lemma}
\newtheorem{remark}[theorem]{Remark}
\newtcolorbox{statusbox}[1][]{colback=yellow!8!white,colframe=orange!70!black,arc=3pt,boxrule=1pt,title=Audit Status,#1}
\title{A2 Audit: Four-Channel Geometric-Mean Measure}
\author{Ing. David Jaroš}
\date{2026-06-14}
\begin{document}
\maketitle
\begin{abstract}
This audit proves the fourth-root insertion used in the G137-B coefficient.  The compact determinant $Z_{\rm compact}$ is common to four equivalent U(1)/Dirac readout channels, $C_Q=4$.  If the channel symmetry is preserved, each channel must receive the same multiplicative weight $w$.  Since the product over four channels must reconstruct the full determinant, $w^4=Z_{\rm compact}$ and hence $w=Z_{\rm compact}^{1/4}$.  Thus the quarter root is not an arbitrary exponent but the unique symmetric per-channel determinant normalisation.
\end{abstract}
\begin{statusbox}
\textbf{A2 result.} Fourth-root measure $Z^{1/4}$ is audited as the unique symmetric per-channel determinant weight for $C_Q=4$.\
\textbf{Residual risk.} Reviewer may challenge exact equivalence of the four channels; that is now the precise audit point.
\end{statusbox}
\section{Claim under audit}
The alpha coefficient uses
\[
  B(\rho)=12^{3/2}Z_{\rm compact}(i\rho)^{1/4}.
\]
The question is why the compact determinant enters with exponent $1/4$.
\section{Channel symmetry}
The information-loss route isolates four equivalent readout channels,
\[
  C_Q=4,
\]
corresponding to the minimal U(1)/Dirac channel package.  The compact determinant is a common factor of the readout sector, not a property of one preferred channel.
\begin{theorem}[Unique symmetric channel weight]
Let $Z$ be a common compact determinant distributed over $C_Q=4$ equivalent readout channels.  If channel permutation symmetry is preserved and the total determinant factor is multiplicative, the unique per-channel weight is
\[
  w=Z^{1/4}.
\]
\end{theorem}
\begin{proof}
By equivalence, all four channels have the same weight $w$.  Multiplicativity of independent readout channels gives
\[
  Z=w\cdot w\cdot w\cdot w=w^4.
\]
For positive determinant weight $Z>0$, the physical branch is $w=Z^{1/4}$.
\end{proof}
\section{Reviewer attack table}
\begin{longtable}{p{0.30\textwidth}p{0.36\textwidth}p{0.24\textwidth}}
\toprule
Attack & Response & Residual risk\\
\midrule\endhead
Why not $Z$ instead of $Z^{1/4}$? & That assigns the whole determinant to one channel and breaks $C_Q=4$ symmetry. & Only if channels are not equivalent.\\
Why not $Z^{1/2}$? & That would imply two equivalent channels, not four. & Requires re-auditing $C_Q$.\\
Why geometric mean? & Determinants multiply; equal multiplicative weights imply geometric mean. & None if multiplicativity accepted.\\
\bottomrule
\end{longtable}
\section{Conclusion}
The exponent $1/4$ follows from channel symmetry and determinant multiplicativity.  A2 is therefore audited internally, conditional on $C_Q=4$ equivalence.
\end{document}
